% Part: first-order-logic
% Chapter: beyond
% Section: second-order-logic

\documentclass[../../../include/open-logic-section]{subfiles}

\begin{document}

\olfileid{fol}{byd}{sol}

\olsection{دوجے درجے دی منطق}

دوجے درجے دی منطق دی زبان وچ صرف افراد دے
!!a{domain} اُتے ہی نہیں، ایس !!{domain} اُتے
بنے تعلقاں اُتے وی مسوّریت ہو سکدی اے۔ پہلے درجے
دی دتی ہوئی زبان $\Lang{L}$ لئی، ہر $k$ واسطے
ایدے وچ !!{variable}s $R$ شامل کیتے جاندے نیں
جیہڑے $k$-جگہی تعلقاں اُتے چل دے نیں، تے انہاں
متغیراں اُتے مسوّریت دی اجازت ہوندی اے۔ جے $R$
اک $k$-جگہی تعلق دا !!a{variable} ہوئے، تے
$t_1$, \dots,~$t_k$ عام (پہلے درجے دیاں)
اصطلاحاں ہون، تاں
$\Atom{R}{t_1,\dots,t_k}$ اک جوہری !!{formula}
اے۔ ہور !!{formula}s دا مجموعہ پہلے درجے دی
منطق وانگ ہی متعین ہوندا اے، بس دوجے درجے دی
مسوّریت لئی وادھو شقّاں ہوندیاں نیں۔ چیتا رہے
کہ !!{identity} صرف پہلے درجے دیاں اصطلاحاں
لئی اے: جے $R$ تے $S$ اکّو مقامیت~$k$ دے
تعلق-!!{variable}s ہون، تاں
$\eq[R][S]$ نوں ایس گل دا مختصر نشان بنا
سکدے آں:
\[
\lforall[x_1 \dots][\lforall[x_k][(\Atom{R}{x_1, \dots, x_k} \liff
  \Atom{S}{x_1, \dots, x_k})]].
\]

دوجے درجے دی منطق دے قاعدے صرف مسوّراں دے
قاعدیاں نوں نویں دوجے درجے دے متغیراں تک
ودھاندے نیں۔ پر ایتھے ذرا دھیان نال سمجھانا
پیندا اے کہ ایہ متغیر~$\Lang{L}$ دے
!!{predicate}s نال، تے عمومی طور اُتے
~$\Lang{L}$ دے !!{formula}s نال، کیویں
مل کے کم کردے نیں۔ گھٹ توں گھٹ تعلقی متغیر
اصطلاحاں ونجوں گنے جاندے نیں، ایس لئی
ایس صورت دے استنباط ہوندے نیں:
\[
!A(R) \Proves \lexists[R][!A(R)]
\]
پر جے $\Lang{L}$ حساب دی زبان ہوئے جیہدے وچ
اک مستقل تعلقی علامت~$<$ ہوئے، تاں ایہ
استنباط وی درست ہون دی توقع ہوئے گی:
\[
x < y \Proves \lexists[R][\Atom{R}{x,y}]
\]
یا دتے ہوئے !!{formula}~$!A$ لئی:
\[
!A(x_1, \dots, x_k) \Proves \lexists[R][\Atom{R}{x_1,\dots,x_k}]
\]
ہور عام طور اُتے اسیں ایس صورت دے استنباط
دی اجازت چاہ سکدے آں:
\[
\Subst{!A}{\lambd[\vec x][!B(\vec x)]}{R} \Proves
\lexists[R][!A]
\]
جتھے $\Subst{!A}{\lambd[\vec x][!B(\vec x)]}{R}$
دا مطلب اے کہ
$\Obj{R}{t_1,\dots,t_k}$ صورت دے ہر جوہری
!!{formula} نوں،~$!A$ وچ جتھے ایس تعلقی متغیر دی آزاد آمد ہووے،
$!B(t_1, \dots, t_k)$ نال بدل
دتا جاوے۔
% PNB-SEM-029: بدل صرف آزاد تعلقی آمد لئی اے؛ پابند ناں تازہ رکھ کے پیرامیٹر گرفت توں بچایا، مشق وچ آزاد آمد واضح کیتی۔
ایہ بدل پابند متغیراں دے نام ضرورت مطابق تازہ رکھ کے
کیتی جاندی اے، تاں جو فارمولے دے آزاد پیرامیٹر غلطی
نال پابند نہ ہو جان؛ عام اصطلاحاں دی بدل وی ایس
گرفت توں بچن والی شرط دے تابع اے۔ ایہ پچھلا قاعدہ {\em خاصیت توں
تعلق بنان دے خاکے} دے برابر اے، یعنی ایس صورت
دا اک مسلّمہ:
\[
\lexists[R][\lforall[x_1, \dots, x_k][(!A(x_1, \dots, x_k) \liff
\Atom{R}{x_1, \dots, x_k})]],
\]
دوجے درجے دی زبان دے ہر !!{formula}~$!A$
لئی اک، جیہدے وچ $R$ آزاد متغیر نہ ہوئے۔
(مشق: وکھاؤ کہ جے $R$ نوں~$!A$ وچ آزاد آون دی
اجازت دتی جاوے، تاں ایہ خاکہ غیر موافق
ہو جاندا اے!)

منطق دان جدوں ”دوجے درجے دی منطق دے
مسلّمے“ آکھدے نیں، تاں عموماً انہاں دی مراد
پہلے درجے دی منطق دے مسوّر-قاعدیاں نوں
دوجے درجے تک گھٹ توں گھٹ ودھانا، تے نال
خاصیت توں تعلق بنان دا خاکہ شامل کرنا ہوندا
اے۔ پر ایدے مسلّمیاں تے قاعدیاں دے ہور
کمزور ذیلی نظاماں دا مطالعہ وی دلچسپ
ہوندا اے۔ مثال لئی چیتا رہے کہ اپنی پوری
عمومیت وچ ایہ خاکہ \emph{غیر اسنادی}
(impredicative) اے: ایہ
$\Atom{R}{x_1, \dots, x_k}$ تعلق دی موجودگی
دا دعویٰ کرن دی اجازت دیندا اے جیہڑا دوجے
درجے دے مسوّراں والے !!a{formula} نال
”متعین“ ہوندا اے؛ تے ایہ مسوّر ایداں دے
سارے تعلقاں دے سیٹ اُتے چل دے نیں---جیہدے
وچ خود $R$ وی شامل اے!{} ویہویں صدی دے
موڑ دے نیڑے رسل دے پیراڈوکس اُتے اک عام
ردّعمل ایہ سی کہ قصور ایداں دیاں تعریفاں
دا اے، تے ریاضی دیاں بنیاداں بناندے ہوئیاں
انہاں توں بچیا جاوے۔ جے
!!{formula}~$!A$ وچ دوجے درجے دے مسوّراں
دی ورتوں منع کر دتی جاوے، تاں خاصیت
توں تعلق بنان دا اک {\em اسنادی}
(predicative) روپ ملدا اے، جیہڑا کجھ
کمزور اے۔

معنیات دے لحاظ نال دوجے درجے دی
!!{structure} نوں ایس طرح سمجھیا جا سکدا
اے: زبان لئی اک پہلے درجے دی
!!{structure}، تے نال ایہو جہے تعلقاں دا
اک سیٹ جیہڑے ایس دے !!{domain} اُتے بنے
ہون تے جیہناں اُتے دوجے درجے دے مسوّر
چل دے ہون (زیادہ عین طور اُتے، ہر $k$
لئی $k$-جگہی تعلقاں دا اک سیٹ)۔
ظاہر اے، جے خاصیت توں تعلق بنان دا خاکہ
!!{derivation} نظام وچ شامل ہوئے، تاں وادھو
شرط ایہ اے کہ ”دوجے درجے دے حصے“ وچ
اینے تعلق ہون کہ اوہ مسلّمے پورے ہون؛
ورنہ !!{derivation} نظام درست نہیں رہندا!{}
کافی تعلق مہیا کرن دا اک آسان طریقہ ایہ
اے کہ دوجے درجے دے حصے وچ پہلے درجے دے
حصے اُتے بن سکدے \emph{سارے} تعلق رکھے
جان۔ ایسی !!a{structure} نوں
\emph{مکمل} (full) آکھدے نیں، تے اک معنی
وچ ایہی زبان دی ”مقصود !!{structure}“
اے۔ جے توجہ صرف مکمل !!{structure}s تک
محدود رکھیے، تاں ایہنوں دوجے درجے دی
\emph{مکمل} معنیات آکھیا جاندا اے۔
ایس صورت وچ !!{structure} متعین کرنا
اصل وچ پہلے درجے دا حصہ متعین کرنا ہی
اے، کیوں جے دوجے درجے دے حصے دے عناصر
اوتھوں آپے متعین ہو جاندے نیں۔

مختصر ایہ کہ دوجے درجے دی منطق دی گل
کردیاں کجھ ابہام رہندا اے۔ !!{derivation}
نظام دے لحاظ نال دو مراداں ہو سکدیاں نیں:
\begin{enumerate}
\item اک ”گھٹ توں گھٹ“ دوجے درجے دا
  !!{derivation} نظام، جیہدے نال خاصیت
  توں تعلق بنان دے کجھ مسلّمے ہون۔
\item ”معیاری“ دوجے درجے دا
  !!{derivation} نظام، جیہدے وچ
  خاصیت توں تعلق بنان دا پورا خاکہ ہوئے۔
\end{enumerate}
معنیات دے لحاظ نال وی دو مراداں نیں:
\begin{enumerate}
\item ”کمزور“ معنیات، جتھے !!a{structure}
  وچ پہلے درجے دا حصہ تے نال دوجے درجے
  دا اینا وڈا حصہ ہوئے جو خاصیت توں
  تعلق بنان دے مسلّمے پورے کرے۔
\item ”معیاری“ دوجے درجے دی معنیات،
  جتھے صرف مکمل !!{structure}s وچاریاں
  جاندیاں نیں۔
\end{enumerate}
جے منطق دان اپنی مراد والا
!!{derivation} نظام یا معنیات واضح نہ کرن،
تاں عموماً ہر فہرست دی دوجی شے مراد
ہوندی اے۔ ایس معنیات دا فائدہ ایہ اے
کہ، جداں اسیں ویکھاں گے، ایہ بہت ساریاں
عام ریاضیاتی ساختاں دا
تماثل تک یکتا بیان دیندی اے؛
نال ای !!{derivation} نظام کافی طاقتور
تے ایس معنیات لئی درست اے۔ نقصان ایہ اے
کہ !!{derivation} نظام ایس معنیات لئی
\emph{تمام} نہیں؛ اصل وچ مکمل دوجے درجے
دی معنیات لئی \emph{کوئی وی} مؤثر طریقے
نال دتا !!{derivation} نظام تمام نہیں۔
دوجے پاسے، اسیں ویکھاں گے کہ
!!{derivation} نظام کمزور معنیات لئی
\emph{تمام} اے؛ یعنی جے کوئی جملہ
ثابت نہیں ہو سکدا، تاں \emph{کوئی نہ
کوئی} !!{structure}، ضروری نہیں کہ
مکمل ہوئے، ایسی اے جیہدے وچ اوہ
جھوٹا اے۔

دوجے درجے دی منطق دی زبان کافی
اظہار-طاقت رکھدی اے۔ اک-جگہی تعلقاں
نوں !!{domain} دے ذیلی سیٹاں نال
پہچانیا جا سکدا اے؛ ایس لئی انہاں
سیٹاں اُتے وی مسوّریت ہو سکدی اے۔
مثال لئی قدرتی اعداد دا استقرا اک
ہی مسلّمے نال بیان کیتا جا سکدا اے:
\[
\lforall[R][((\Atom{R}{\Obj{0}} \land \lforall[x][(\Atom{R}{x} \lif
    \Atom{R}{x'})]) \lif \lforall[x][\Atom{R}{x}])].
\]
جے حساب دی زبان دیاں علامتاں
$\Obj 0, \Obj \prime, +,
\times$ تے $<$ ہون، تاں انہاں دے
رویے لئی ہیٹھلے مسلّمے شامل
کیتے جا سکدے نیں:
\begin{enumerate}
\item $\lforall[x][\lnot x' = \Obj 0]$
% PNB-SYN-010: اصل دا غیر مقرر s اوپر دتی جانشین علامت دی بجائے آ گیا سی؛ اوہی پرائم علامت ورتی۔
\item $\lforall[x][\lforall[y][(x' = y' \lif x = y)]]$
\item $\lforall[x][(x + \Obj 0) = x]$
\item $\lforall[x][\lforall[y][(x + y') = (x + y)']]$
\item $\lforall[x][(x \times \Obj 0) = \Obj 0]$
\item $\lforall[x][\lforall[y][(x \times y') = ((x \times y) + x)]]$
\item $\lforall[x][\lforall[y][(x < y \liff \lexists[z][y = (x + z')])]]$
\end{enumerate}
ایہ وکھانا اوکھا نہیں کہ ایہ مسلّمے،
اوپرلے استقرا دے مسلّمے نال مل کے،
حساب دے معیاری نمونے
!!{structure}~$\Struct{N}$ دا
تماثل تک یکتا بیان دیندے نیں،
بشرطیکہ مکمل دوجے درجے دی معنیات
ورتی جا رہی ہوئے۔ ایداں دی کوئی وی
!!{structure}~$\Struct{M}$ لو جیہدے
وچ ایہ مسلّمے سچے نیں۔ اک
فنکشن~$f$، $\Nat$ توں
$\Struct{M}$ دے !!{domain} تک،
$\Nat$ اُتے عام بازگشت نال
ایس طرح متعین کرو کہ
$f(0) = \Assign{\Obj 0}{M}$ تے
$f(x+1) = \Assign{\prime}{M}(f(x))$
ہوئے۔ $\Nat$ اُتے عام استقرا نال،
تے ایس گل نال کہ مسلّمے (1)
تے~(2)~$\Struct M$ وچ پورے ہوندے
نیں، ویکھدے آں کہ $f$ !!{injective}
اے۔ $f$ دے !!{surjective} ہون
لئی~$P$ نوں~$\Domain{M}$ دے انہاں
عناصر دا سیٹ لو جیہڑے~$f$ دی
حاصل قدراں وچ نیں۔ چونکہ
$\Struct M$ مکمل اے، $P$ دوجے درجے
دے !!{domain} وچ شامل اے۔
~$f$ دی بناوٹ توں پتا اے کہ
$\Assign{\Obj
  0}{M}$ سیٹ~$P$ وچ اے، تے
$P$،~$\Assign{\prime}{M}$ دے
عمل ہیٹھ بند اے۔ استقرا دا مسلّمہ
$\Struct M$ وچ (خاص طور اُتے~$P$
لئی) پورا ہون دی وجہ توں
$P$~پہلے درجے دے پورے !!{domain}
دے برابر اے، جیہڑا~$\Struct M$ دا
دائرہ اے۔ ایس نال $f$ دا
!!a{bijection} ہون ثابت ہوندا اے۔
$f$ دا ہومومارفزم ہون وکھانا وی
اوداں ای آسان اے: $\Nat$ اُتے
عام استقرا وار وار ورتو۔

سیٹ نظریے دی زبان وچ فنکشن تعلق
دی اک خاص قسم ہوندا اے؛ مثال لئی
اک-جگہی فنکشن~$f$ نوں ایداں دے
دو-جگہی تعلق~$R$ نال پہچانیا
جا سکدا اے جیہڑا
$\lforall[x][\lexists![y][R(x,y)]]$
پورا کرے۔ ایس لئی فنکشناں اُتے
وی مسوّریت ہو سکدی اے۔ مکمل
معنیات ورت کے لامتناہی
!!{structure}s دی جماعت نوں
انہاں !!{structure}s $\Struct M$
دی جماعت ونجوں متعین کیتا جا سکدا
اے جیہناں لئی $\Struct M$ دے
!!{domain} توں ایدے اپنے کسے
حقیقی ذیلی سیٹ ول اک
انجیکٹو فنکشن ہوئے:
\[
\lexists[f][(\lforall[x][\lforall[y][(\eq[f(x)][f(y)] \lif \eq[x][y])]]
  \land \lexists[y][\lforall[x][\eq/[f(x)][y]]])].
\]
ایس جملے دی نفی فیر متناہی
!!{structure}s دی جماعت متعین
کردی اے۔
% PNB-SEM-030: اپنے حقیقی ذیلی سیٹ نال برابر سائز دا بیان ہر لامتناہی دائرے لئی آکھن وچ معمول دا انتخابی مفروضہ ورتیا گیا اے۔
ایتھے ہر لامتناہی سیٹ وچ اک گننے جوگا لامتناہی
ذیلی سیٹ ہون دا معمول دا مفروضہ وی ورتیا گیا اے،
مثال لئی انتخاب دے مسلّمے والے سیٹ نظریاتی پس منظر
وچ۔ ایس مفروضے توں بغیر اوپرلا فارمولا بالخصوص
اپنے حقیقی ذیلی سیٹ نال برابر سائز رکھن والے دائرے
بیان کردا اے؛ ایہدی نفی نوں صرف متناہی دائرے ہون
دا عین بیان نہیں آکھ سکدے۔

ایس توں علاوہ، خطی ترتیب دی
تعریف وچ ہیٹھلا اضافہ کرکے
حُسنِ ترتیب والیاں ساختاں دی
جماعت متعین کیتی جا سکدی اے:
\[
\lforall[P][(\lexists[x][\Atom{P}{x}] \lif \lexists[x][(\Atom{P}{x}
    \land \lforall[y][(y < x \lif \lnot \Atom{P}{y})])])].
\]
ایہ آکھدا اے کہ ہر غیر خالی
سیٹ دا سب توں چھوٹا عنصر ہوندا
اے، جے ”سیٹ“ نوں ”اک-جگہی
تعلق“ نال اک من لیا جاوے۔ اک
ہور مثال وچ گرافاں دی جڑت دا
تصور ایس طرح بیان ہو سکدا اے
کہ راساں نوں دو غیر خالی
حصیاں وچ ایداں وکھ نہیں
کیتا جا سکدا کہ انہاں
وچکار کوئی کنارہ نہ ہوئے۔ ایتھے گراف دے کنارے
بغیر سمت دے منے گئے نیں، یعنی دو راساں دا تعلق
دوواں پاسیاں توں اکّو اے:
% PNB-SEM-031: گراف دی عام جڑت لئی بغیر سمت والا تعلق واضح کیتا؛ سمت والے گراف دی کمزور جڑت دا دعویٰ نہیں۔
\[
\lnot \lexists[A][(\lexists[x][A(x)] \land \lexists[y][\lnot A(y)]
  \land \lforall[w][\lforall[z][((\Atom{A}{w} \land \lnot \Atom{A}{z})
      \lif \lnot \Atom{R}{w,z})]])].
\]
ہور مشق ونجوں تسی انہاں متناہی
!!{structure}s دی جماعت متعین کرن
دی کوشش کر سکدے او جیہناں دے
!!{domain} دا سائز جفت اے۔
ایدے توں وی ودھ نمایاں گل ایہ
اے کہ ناطق اعداد نوں شامل کرن
والے مکمل ترتیب دار میدان ونجوں
حقیقی اعداد دا
تماثل تک یکتا بیان دتا
جا سکدا اے۔

مختصر ایہ کہ دوجے درجے دی
منطق پہلے درجے دی منطق نالوں
کیتے ودھ اظہار-طاقت رکھدی
اے۔ ایہ چنگی خبر اے؛ ہن بری
خبر: اسیں دس چکے آں کہ
مکمل دوجے درجے دی معنیات
لئی کوئی مؤثر !!{derivation}
نظام تمام نہیں۔ چاہے ایہ
چنگا لگے یا برا، پہلے درجے
دی منطق دیاں کئی خاصیتاں
ایتھے نہیں ملتیاں، جنہاں وچ
محدود گواہی (compactness) تے
L\"owenheim--Skolem دے
قضیے وی شامل نیں۔

دوجے پاسے، جے کوئی مکمل
دوجے درجے دی معنیات چھڈ
کے کمزور معنیات منّن نوں
تیار ہوئے، تاں گھٹ توں
گھٹ دوجے درجے دا
!!{derivation} نظام ایس
معنیات لئی تمام اے۔
دوجے لفظاں وچ جے
$\Proves$ نوں ”گھٹ توں گھٹ
نظام وچ ثابت کردا اے“
تے $\Entails$ نوں
”کمزور معنیات وچ منطقی
نتیجہ اے“ پڑھیا جاوے،
تاں جدوں وی
$\Gamma \Entails !A$،
اوہدوں
$\Gamma \Proves !A$
وکھایا جا سکدا اے۔ جے
!!{derivation} نظام وچ
خاص خاص خاصیت توں تعلق
بنان دے مسلّمے شامل
کرنے ہون، تاں معنیات
نوں انہاں مسلّمیاں پورے
کرن والیاں دوجے درجے
دیاں ساختاں تک محدود
کرنا پیندا اے: مثال
لئی جے $\Delta$ ایداں
دے مسلّمیاں دا سیٹ
ہوئے (ممکن اے سارے
ای ہون)، تاں
$\Gamma
\cup \Delta \Entails !A$
توں
$\Gamma \cup \Delta \Proves !A$
ملدا اے۔ خاص طور
اُتے، جے $!A$ انہاں
خاصیت توں تعلق بنان
دے مسلّمیاں نال
ثابت نہ ہو سکدا
ہوئے، تاں
$\lnot !A$ دا اک
نمونہ موجود اے
جیہدے وچ ایہ
مسلّمے فیر وی
پورے ہوندے نیں۔

کمزور معنیات لئی
تمامیت دا قضیہ
کیوں اے، ایہ ویکھن
دا سب توں آسان
طریقہ دوجے درجے
دی منطق نوں کئی
قسماں والی منطق
سمجھنا اے۔ اک قسم
عام ”پہلے درجے“
دے !!{domain}
ونجوں سمجھی جاندی
اے، تے فیر ہر $k$
لئی ”$k$-جگہی
تعلقاں“ دا
!!a{domain}
ہوندا اے۔ زبان
وچ پہلے توں
شامل تعلقی
علامتاں
”$\Atom{\Obj{true}_k}{R,x_1,\dots,x_k}$“
رکھدے آں؛ ایس
دا مطلب ایہ
اے کہ $R$،
$x_1$, \dots,~$x_k$
لئی پورا
اتردا اے،
جتھے $R$
”$k$-جگہی تعلق“
قسم دا متغیر
اے تے
$x_1$, \dots,~$x_k$
پہلے درجے
دی قسم دیاں
شےواں نیں۔

ایس پہچان
نال کمزور
دوجے درجے
دی معنیات
اصل وچ کئی
قسماں والی
منطق دی عام
معنیات ہی
اے؛ تے اسیں
پہلاں ویکھ
چکے آں کہ
کئی قسماں
والی منطق
پہلے درجے
دی منطق وچ
سموئی جا
سکدی اے۔
دوواں پاسیاں
دے ترجمیاں
نوں پیش نظر
رکھ کے، دوجے
درجے دی منطق
دا کمزور تصور
اصل وچ پہلے
درجے دی منطق
دا ہی اک
بھیس بدلیا
روپ اے،
جتھے
!!{domain}
وچ ”شےواں“
تے ”تعلق“
دوویں ہوندے
نیں تے مناسب
مسلّمے انہاں
اُتے لاگو
ہوندے نیں۔

\end{document}
